\section{The real estimate: repulsion balanced against decay}
\label{analytic:section}

The determinant is positive, but positivity alone says nothing about its
size. Its integral representation contains two competing effects. The
Vandermonde factor favours separated integration points; the weight
penalizes points far from the origin. We now bound that competition
uniformly, including configurations in which points nearly collide.
The method, the comparison measure, and its numerical data are due to
Fauzan \cite[Section 6 and Appendix A]{fauzan}.
Firsching's formalization supplies an additional detailed implementation
of these arguments \cite[\texttt{Apery/Energy*.lean},
\texttt{Apery/RealBound.lean}]{firsching}.
The proofs below spell out the real-variable argument and independently
recheck the finite numerical certificate with rational enclosures.

Throughout this section,
\[
 \alpha=\frac3{40},\qquad \lambda=\frac{37}{40},\qquad
 N=\alpha K,\qquad h=\lambda K.
\]
Recall that \(F_K=S_K\Delta_K\). We shall prove
\begin{equation}\label{analytic:goal}
 0<F_K(\zeta(5)),\qquad
 \log F_K(\zeta(5))
 \leq -\frac{2733991}{2000000}K^2+27K\log K+200K.
\end{equation}
All logarithms in this section are natural.

\subsection{From a Gram determinant to a many-variable integral}

Recall \(d\nu_K(y)=A_K(y^2)w(y)\,dy\), where
\(A_K=D_N^6/D_K\). The already established moment identity gives
\(G_K(\zeta(5))_{ij}=\int y^{2i+2j}\,d\nu_K(y)\).
For any integrable products \(f_i g_j\), the relevant form of
Andr\'eief's identity \cite{andreief} is
\[
 \int \det[f_i(y_j)]_{i,j=0}^{h-1}
       \det[g_i(y_j)]_{i,j=0}^{h-1}
       \prod_{j=0}^{h-1}d\nu_K(y_j)
 =h!\det\left[\int f_i g_j\,d\nu_K\right]_{i,j=0}^{h-1}.
\]
Here is the proof in the situation we need. Expand both determinants as
finite sums over permutations. Fubini separates each product into \(h\)
one-variable integrals. Relabeling the integration variables by the
permutation in the first determinant leaves \(h!\) identical copies of
the determinant expansion on the right. All these interchanges are
legitimate: \(A_K(y^2)\) has polynomial growth and \(w(y)\) has exponential
decay, so every polynomial moment occurring in the finite expansion is
absolutely integrable. Taking \(f_i=g_i=y^{2i}\) proves
\begin{equation}\label{analytic:andreief}
 \Delta_K(\zeta(5))=
 \frac1{h!}\int_{(0,\infty)^h}
 \prod_{i<j}(y_i^2-y_j^2)^2
 \prod_{i=1}^h A_K(y_i^2)w(y_i)\,dy_1\cdots dy_h .
\end{equation}
This is the classical determinant identity used in
\cite[(6.10)]{fauzan}; its application here is also formalized in
\cite[\texttt{Apery/Andreief.lean},
\texttt{Apery/DetIntegral.lean}]{firsching}.

The natural scale is \(y_i=K\sqrt{t_i}\). The finite products explain
the resulting field:
\[
 \log D_m(K^2t)=2m\log K+
                   \sum_{j=1}^m\log\bigl(t+(j/K)^2\bigr).
\]
The sum differs from \(K\int_0^{m/K}\log(t+u^2)\,du\) by
\(O(\log K)\), uniformly for \(t\geq0\) and \(1\leq m\leq K\);
\eqref{analytic:riemann} proves the explicit bound, including \(t=0\).
Applying this to \(D_N^6/D_K\) and including the exponential decay
of \(w\) produces the external field
\begin{equation}\label{analytic:field}
 V(t)=2\pi\sqrt t+\int_0^1\log(t+u^2)\,du
                   -6\int_0^\alpha\log(t+u^2)\,du,\qquad t\geq0.
\end{equation}
The term \(2\pi\sqrt t\) comes from the exponential decay of \(w\); the
two logarithmic integrals approximate the products in \(D_N^6/D_K\).
The principal exponent in the rescaled integrand is
\[
 2\sum_{i<j}\log|t_i-t_j|-K\sum_i V(t_i).
\]
We first bound this expression. Afterwards we return to all the factors
discarded in identifying its principal part.

The same product data reappear arithmetically: prime valuations give
divisibility counts, with Legendre's floors for factorials.
Each estimate needs its own error bounds
\cite[Sections 5--6]{fauzan}.

\subsection{An explicit comparison measure}
\label{analytic:measure}

For a finite measure \(\rho\), write
\[
 U^\rho(t)=\int\log|t-u|\,d\rho(u),\qquad
 I(\rho)=\iint\log|t-u|\,d\rho(t)\,d\rho(u).
\]
The sign convention matters: the logarithm appears with a positive sign.
We need a measure of mass \(\lambda\) satisfying
\begin{equation}\label{analytic:potential-target}
 2U^\rho(t)-V(t)\leq M_0,\qquad
 M_0=-\frac{1329}{200}=-6.645,\qquad t\geq0.
\end{equation}
It is enough to exhibit one suitable measure. No assertion that it is
the exact equilibrium measure is needed.

To see why, let \(\eta,\rho\) be positive, compactly supported measures
on \([0,\infty)\), both of mass \(\lambda\), with \(\rho\) satisfying
\eqref{analytic:potential-target} and
\(\lvert\log|t-u|\rvert\) integrable against
\((\eta+\rho)\otimes(\eta+\rho)\).
Their difference has mass zero, so Lemma~\ref{analytic:zero-mass}
below gives \(I(\eta-\rho)\leq0\). Expanding yields
\begin{align*}
 I(\eta)-\int V\,d\eta
 &=I(\eta-\rho)+\int(2U^\rho-V)\,d\eta-I(\rho)\\
 &\leq\lambda M_0-I(\rho).
\end{align*}
Thus one explicit potential bounds every admissible configuration's
energy \cite[Section 6.1]{fauzan}.

For actual points, \(K^{-1}\sum_i\delta_{t_i}\) has the required mass
\(h/K=\lambda\), but infinite negative self-energy. Spreading its
atoms over small circles gives the measure \(\sigma\) used below.
The three roles are distinct: \(\nu_K\) defines the determinant's
moments, \(\sigma\) describes one smoothed configuration, and \(\rho\)
certifies a uniform energy bound. Fauzan chooses the latter for its
tractable potential \cite[after (6.4)]{fauzan}.

Fauzan's measure is
\[
 \rho=\sum_{j=1}^{16}c_j\omega_{[a_j,b_j]},\qquad
 d\omega_{[a,b]}(u)=\frac{\mathbf 1_{(a,b)}(u)\,du}
                         {\pi\sqrt{(u-a)(b-u)}}.
\]
The following integer table specifies the data exactly: each entry in
the last three columns is divided by \(10^{12}\).
The data are reproduced from \cite[Table 1]{fauzan}; the transcription
was checked against \cite[\texttt{Apery/Measure.lean},
\texttt{scripts/lemma61.py}]{firsching}.

\begin{center}
\small
\begin{tabular}{r r r r}
\(j\)&\(10^{12}a_j\)&\(10^{12}b_j\)&\(10^{12}c_j\)\\ \hline
1&3906748086&8992695531&10515596180\\
2&2312248264&15340997855&29471737793\\
3&1402286665&25730180724&42934365099\\
4&881725356&41909578246&58204231966\\
5&578197906&65851089563&69037621310\\
6&396324613&99481037884&78873099189\\
7&283911191&144325727458&84856120711\\
8&212206188&201105762729&88396082127\\
9&165097686&269345996903&88303382125\\
10&133347132&347089554156&85472321255\\
11&111522114&430806704415&78899184238\\
12&96349355&515561896511&70353471918\\
13&85815639&595448778546&58838976615\\
14&78667711&664241383483&44421321106\\
15&74129565&716160577112&30462865791\\
16&71741310&746637295669&5959622577
\end{tabular}
\end{center}

Integer comparisons and addition give
\begin{equation}\label{analytic:data-facts}
 \sum_jc_j=\lambda,\qquad
 0<a_{16}<\cdots<a_1<b_1<\cdots<b_{16}<2,\qquad
 b_j-a_j>\frac1{225}.
\end{equation}
Thus the intervals are nested, with the first interval innermost.

We shall use two elementary potential formulas. Normalized arclength
\(\omega_{z_0,\varepsilon}\) on a circle satisfies
\begin{equation}\label{analytic:circle}
 \int\log|z-u|\,d\omega_{z_0,\varepsilon}(z)
       =\log\max(|z_0-u|,\varepsilon).
\end{equation}
For unequal radii this follows by expanding
\(\log|1-re^{i\theta}|\) for \(0\leq r<1\) and integrating its cosine
series: every nonconstant term has mean zero. At equality of the radii,
the remaining identity is the integrable formula
\(\int_0^{2\pi}\log|1-e^{i\theta}|\,d\theta=0\).
For example, the latter follows from
\(\int_0^{\pi/2}\log\sin\theta\,d\theta
=-\frac\pi2\log2\), obtained by the substitutions
\(\theta\mapsto\pi/2-\theta\) and \(u=2\theta\).

For a unit arcsine measure, with \(m=(a+b)/2\) and \(r=(b-a)/2\),
\begin{equation}\label{analytic:arc-potential}
 U^{\omega_{[a,b]}}(t)=
 \begin{cases}
 \log((b-a)/4),&a\leq t\leq b,\\[2mm]
 \displaystyle\log\frac{|t-m|+\sqrt{(t-a)(t-b)}}2,
     &t\notin[a,b].
 \end{cases}
\end{equation}
One can obtain this directly from \eqref{analytic:circle}. The change
\(u=m+r\cos\theta\) turns arcsine averaging into uniform circle
averaging. With \(z=e^{i\theta}\),
\[
 t-m-\frac r2(z+z^{-1})
   =-\frac r{2z}(z-q_+)(z-q_-),\qquad
 q_\pm=\frac{t-m}{r}\pm\sqrt{\left(\frac{t-m}{r}\right)^2-1}.
\]
The two roots have product one. Inside the interval both have modulus
one; outside it exactly one has modulus greater than one. Applying the
circle formula to these two factors gives
\eqref{analytic:arc-potential}, including its continuous endpoint values.
This is the classical arcsine calculation used in
\cite[(A.1)]{fauzan}.

Let \(s_j=\sum_{i\leq j}c_i\), with \(s_0=0\).
The potential of the \(j\)-th interval is constant on every earlier
interval. Therefore both the self-energy and the mutual energies can
be evaluated without quadrature:
\begin{equation}\label{analytic:energy}
 I(\rho)=\sum_{j=1}^{16}(s_j^2-s_{j-1}^2)
                       \log\frac{b_j-a_j}{4}.
\end{equation}
Indeed, the coefficient is
\(c_j^2+2c_j\sum_{i<j}c_i=s_j^2-s_{j-1}^2\).
This is exactly the use of nesting in \cite[(A.2)]{fauzan}.

\subsection{Verifying the potential bound everywhere}
\label{analytic:certificate}

The remaining task is a uniform inequality, not an estimate at a
collection of sample points. We give the reduction to finitely many
rational comparisons, together with the supplied certificate.
The partition and reduction below are Fauzan's
\cite[Appendix A.2--A.3, Table 2]{fauzan};
the companion verification program is a separate implementation using
outward-rounded rational intervals. The corresponding Lean certificates
are in \cite[\texttt{Apery/TableCheck.lean},
\texttt{Apery/Table/}]{firsching}.

For \(t>0\), elementary integration in \eqref{analytic:field} yields
\begin{align}
 V(t)={}&\log(1+t)-6\alpha\log(t+\alpha^2)-2+12\alpha \notag\\
       &+2\sqrt t\left(\pi+\arctan\frac1{\sqrt t}
                               -6\arctan\frac{\alpha}{\sqrt t}\right).
 \label{analytic:field-closed}
\end{align}
At \(t=0\), its continuous value is
\(-12\alpha\log\alpha-2+12\alpha\).
Writing \(\Phi(y)=V(y^2)\), differentiation gives
\[
 \Phi'(y)=2\left(\pi+\arctan(1/y)-6\arctan(\alpha/y)\right),
 \qquad
 \Phi''(y)=-\frac2{1+y^2}+\frac{12\alpha}{\alpha^2+y^2}.
\]
The second derivative is positive before \(y^2=711/880\) and negative
afterwards. Since \(\Phi'(0+)=-3\pi\) and
\(\Phi'(+\infty)=2\pi>0\), the derivative has exactly one zero.
Consequently \(V\) decreases to a single minimum and then increases.

Set
\[
 q_-=\frac{59205077}{10^{10}},\qquad
 q_+=\frac{59205079}{10^{10}}.
\]
The rational series enclosures described below give
\(\Phi'(\sqrt{q_-})<0<\Phi'(\sqrt{q_+})\).
Thus the minimum is between these two rational numbers. On that
short interval a lower bound for \(V\) is
\begin{align}
 V_*={}&\log(1+q_-)-6\alpha\log(q_++\alpha^2)-2+12\alpha\notag\\
 &+2\sqrt{q_+}\left(\pi+\arctan\frac1{\sqrt{q_+}}
                         -6\arctan\frac{\alpha}{\sqrt{q_-}}\right).
 \label{analytic:vstar}
\end{align}
The parenthesis is negative, as the same certificate verifies.
Each monotone factor in \eqref{analytic:field-closed} has therefore
been replaced in the direction needed for a lower bound. In particular,
using the larger square root in front of the negative parenthesis is
essential.

The arcsine formula shows that \(U^\rho\) is nonincreasing on
\([0,a_1]\), constant on \([a_1,b_1]\), and nondecreasing afterwards.
Its maximum on any interval is therefore attained at an endpoint.
For \(0\leq l<r\leq2\), put
\[
 B(l,r)=2\max(U^\rho(l),U^\rho(r))-
 \begin{cases}
 V(r),&r\leq q_-,\\
 V(l),&l\geq q_+,\\
 V_*,&\text{otherwise}.
 \end{cases}
\]
The two monotonicity statements prove
\[
 2U^\rho(t)-V(t)\leq B(l,r)\qquad(l\leq t\leq r).
\]
The finite calculation concerns these upper bounds.

To specify the partition, define
\[
\begin{gathered}
 A_0=0,\qquad A_j=a_{17-j}\quad(1\leq j\leq16),\\
 A_{17}=q_-,\qquad A_{18}=q_+,\qquad A_{35}=2,\\
 A_{18+j}=b_j\quad(1\leq j\leq16).
\end{gathered}
\]
Each triple \((j,d,k)\) in the accompanying data file denotes
\begin{equation}\label{analytic:partition}
 \left[A_j+(A_{j+1}-A_j)\frac{k}{2^d},
       A_j+(A_{j+1}-A_j)\frac{k+1}{2^d}\right].
\end{equation}
The complete partition data appear in
\texttt{checks/\allowbreak analytic\_data.json}.
This is the 91-row compressed list from Fauzan's Table 2. Expanding its integer
ranges gives 684 intervals. The verifier sorts their rational endpoints
and checks that the first is zero, the last is two, and every right
endpoint equals the next left endpoint exactly. Thus no part of
\([0,2]\) is omitted.

For clarity, the elementary enclosures used by the verifier are as
follows. If \(1\leq x\leq2\) and \(z=(x-1)/(x+1)\), then
\begin{equation}\label{analytic:log-series}
 0\leq\log x-2\sum_{j=0}^{m-1}\frac{z^{2j+1}}{2j+1}
 \leq\frac{2z^{2m+1}}{(2m+1)(1-z^2)}.
\end{equation}
This follows by integrating the geometric series for \(2/(1-z^2)\);
bounding the denominators in its positive tail gives the displayed
remainder. A general positive argument is reduced by
\(\log(2^e x)=e\log2+\log x\).
For \(0\leq x\leq1/2\), the alternating series gives
\begin{equation}\label{analytic:atan-series}
 \sum_{j=0}^{2m-1}\frac{(-1)^j x^{2j+1}}{2j+1}
 \leq \arctan x
 \leq \sum_{j=0}^{2m-1}\frac{(-1)^j x^{2j+1}}{2j+1}
                    +\frac{x^{4m+1}}{4m+1}.
\end{equation}
Larger arguments are reduced using
\[
 \arctan x=\frac\pi2-\arctan(1/x)\ (x>1),\qquad
 \arctan x=2\arctan\frac{x}{1+\sqrt{1+x^2}}\ (0\leq x\leq1).
\]
Finally \(\pi=16\arctan(1/5)-4\arctan(1/239)\); the tangent addition
formula and the ranges of the angles verify this identity.
Square-root bounds are obtained by integer square roots: if
\(a\leq x\leq b\), integers bracketing
\(\sqrt{a\,2^{224}}\) and \(\sqrt{b\,2^{224}}\) give bounds on
the grid \(2^{-112}\).

The supplied \texttt{analytic\_certificate.py} uses 48 terms in
\eqref{analytic:log-series} and 64 terms in the alternating sum in
\eqref{analytic:atan-series}. Every addition, multiplication, division,
and square root is rounded outwards to that dyadic grid. The remainder
bounds are added before any comparison is accepted. Thus rounding
cannot turn an unproved inequality into a proved one.
The complete list of interval bounds is written to
\texttt{analytic\_intervals.csv}; the exact rational endpoints and file
hashes are recorded in \texttt{analytic\_certificate.json}.
The checks give, for every interval,
\begin{equation}\label{analytic:finite-result}
 B(l,r)<-\frac{6645002}{10^6}<-\,\frac{1329}{200}.
\end{equation}
The largest upper bound occurs at \((j,d,k)=(30,6,26)\) and is less than
\(-6.645002689215588323\). This last decimal is only a readable report
of the rational comparison; the comparisons themselves use integers.
Notice the distinction between the strict certificate target
\(-6.645002\) and the constant \(M_0=-6.645\) used in the proof.

It remains to cover \(t\geq2\). Since \(\rho\) is supported in \((0,2)\),
\(U^\rho(t)\leq\lambda\log t\). Also
\[
 \int_0^1\log(t+u^2)\,du\geq\log t,\qquad
 \int_0^\alpha\log(t+u^2)\,du
       \leq\alpha\log t+\frac{\alpha^3}{t};
\]
the latter follows from \(\log(1+v)\leq v\) and \(u^2\leq\alpha^2\).
Consequently
\begin{equation}\label{analytic:tail}
 2U^\rho(t)-V(t)
 \leq \frac{13}{10}\log t+\frac{6\alpha^3}{t}-2\pi\sqrt t.
\end{equation}
The right-hand side is decreasing for \(t\geq2\), and its value at two
is smaller than \(M_0\). For example, the inequalities
\(\log2<7/10,\ \pi>3,\ \sqrt2>7/5\) already suffice.
Together with \eqref{analytic:finite-result}, this proves
\eqref{analytic:potential-target} on the entire half-line.

\subsection{Why the comparison controls every configuration}

We include the energy argument because it explains both the comparison
method and the treatment of collisions.
The proof follows \cite[Lemma 6.2 and Section 6.1]{fauzan};
the Gaussian-kernel and regularization steps are separately formalized
in \cite[\texttt{Apery/Gaussian.lean},
\texttt{Apery/ZeroMass.lean}, \texttt{Apery/EnergyLimit.lean},
\texttt{Apery/EnergyFinal.lean}]{firsching}.

\begin{lemma}[Zero-mass logarithmic energy]\label{analytic:zero-mass}
Let \(\eta\) be a finite real signed measure with compact support in
\(\mathbb C\), total mass zero, and
\(\iint|\log|z-u||\,d|\eta|(z)\,d|\eta|(u)<\infty\).
Then \(I(\eta)\leq0\).
\end{lemma}
\begin{proof}
For \(s>0\), completing the square in a planar Gaussian integral gives
\[
 e^{-s|z-u|^2}=\frac{4s}{\pi}
 \int_{\mathbb C}e^{-2s|z-v|^2}e^{-2s|u-v|^2}\,dA(v).
\]
Integrating against \(d\eta(z)d\eta(u)\) yields
\[
 J(s):=\iint e^{-s|z-u|^2}\,d\eta(z)d\eta(u)
 =\frac{4s}{\pi}\int_{\mathbb C}
       \left(\int e^{-2s|z-v|^2}\,d\eta(z)\right)^2dA(v)\geq0.
\]
Fubini is valid because the Gaussian convolution is bounded and the
total variation is finite.
For \(0<a<b<\infty\), put
\[
 L_{a,b}(r)=\frac12\int_a^b\frac{e^{-s}-e^{-sr^2}}s\,ds.
\]
Writing \(e^{-s}-e^{-sr^2}=\int_1^{r^2}s e^{-sv}\,dv\)
shows that \(L_{a,b}(r)\to\log r\) as \(a\downarrow0,b\uparrow\infty\).
The integrand has one sign for each fixed \(r\), so
\(|L_{a,b}(r)|\leq|\log r|\).
The mass-zero condition removes the constant \(e^{-s}\) term:
\[
 \iint L_{a,b}(|z-u|)\,d\eta(z)d\eta(u)
       =-\frac12\int_a^b\frac{J(s)}s\,ds\leq0.
\]
The assumed absolute logarithmic integrability makes the diagonal a
null set for \(|\eta|\otimes|\eta|\). Dominated convergence therefore
gives the result.
\end{proof}

For real points \(t_1,\ldots,t_h\geq0\), replace each point mass by
normalized arclength \(\omega_i\) on the circle of radius
\(\varepsilon\) centred at \(t_i\), and set
\(\sigma=K^{-1}\sum_i\omega_i\).
Both \(\sigma\) and \(\rho\) have mass \(\lambda\).
The potential bound is on the real axis; using complex circles requires
the circle-average error estimate below.
We verify explicitly that the lemma applies to \(\sigma-\rho\).
The circle potential \eqref{analytic:circle} is finite and continuous.
For an arcsine measure on an interval of length \(L\), the elementary
inequality
\[
 \frac1{\sqrt{x(L-x)}}\leq
 \frac1{\sqrt L}\left(\frac1{\sqrt x}+\frac1{\sqrt{L-x}}\right)
 \qquad(0<x<L)
\]
shows that any interval of length \(d\) has mass at most
\(4\sqrt d/(\pi\sqrt L)\). In particular, the mass within distance \(r\)
of any point of \(\mathbb C\) is at most
\(4\sqrt{2r}/(\pi\sqrt L)\).
Integrating this estimate against \(dr/r\) at zero proves uniform
integrability of the logarithmic singularity. It also proves continuity
of the arcsine potential, by first truncating the logarithm and then
letting the truncation radius tend to zero.
Thus the potential of the positive measure \(\sigma+\rho\) is bounded
on its compact support. The positive part of the logarithmic kernel
is bounded there as well. Its negative part consequently has finite
double integral. Since
\(|\sigma-\rho|\leq\sigma+\rho\), the hypothesis of
Lemma~\ref{analytic:zero-mass} holds.

The same mass estimate controls the effect of smoothing.
By \eqref{analytic:circle}, Tonelli, and
\(\log(\varepsilon/d)=\int_d^\varepsilon dr/r\),
\begin{align*}
 \int U^\rho\,d\omega_i-U^\rho(t_i)
 &=\int_{|u-t_i|<\varepsilon}
                  \log\frac{\varepsilon}{|u-t_i|}\,d\rho(u)\\
 &\leq \sum_j\frac{8c_j\sqrt{2\varepsilon}}
                         {\pi\sqrt{b_j-a_j}}
 \leq60\sqrt\varepsilon.
\end{align*}
The final inequality follows from \eqref{analytic:data-facts},
\(\sum c_j=37/40\), and \(\pi>3,\sqrt2<3/2\).
The circle identity gives
\[
 I(\omega_i)=\log\varepsilon,\qquad
 \iint\log|z-u|\,d\omega_i(z)d\omega_j(u)
       \geq\log|t_i-t_j|.
\]
For the second inequality, average first over one circle, discard the
maximum with \(\varepsilon\), and then average over the other.

Expanding \(I(\sigma-\rho)\leq0\), multiplying by \(K^2\), and using
these estimates now gives
\begin{align}
 2\sum_{i<j}\log|t_i-t_j|-K\sum_iV(t_i)
 \leq{}&(\lambda M_0-I(\rho))K^2\notag\\
       &+120Kh\sqrt\varepsilon-h\log\varepsilon.
 \label{analytic:configuration-raw}
\end{align}
To see the factors of \(K\), the cross term is
\(2K\sum_i\int U^\rho\,d\omega_i\).
After smoothing and subtracting the field term, it is bounded by
\[
 K\sum_i(2U^\rho(t_i)-V(t_i))+120Kh\sqrt\varepsilon.
\]
The potential inequality bounds its first part by \(KhM_0\).
Coincident points are harmless: the left side is then \(-\infty\).

A constant bound on an unbounded integration domain would not suffice.
We retain a decaying factor by applying the same calculation to the
field \(V(t)-\sqrt t/K\).
For \(t\leq2\), \eqref{analytic:potential-target} immediately gives
\[
 2U^\rho(t)-V(t)+\frac{\sqrt t}{K}
       \leq M_0+\frac{\sqrt2}{K}.
\]
For \(t\geq2\) and \(K\geq2\), use \eqref{analytic:tail} with
\(2\pi\) replaced by \(2\pi-\frac12\).
Its derivative is negative on this range, and its value at two is
smaller than \(M_0\): the same elementary bounds give
\[
 \frac{13}{10}\log2+3\alpha^3
       -(2\pi-\tfrac12)\sqrt2
 <0.91+3(3/40)^3-5.5(1.4)<-6.645.
\]
Thus the displayed bound holds for all \(t\geq0\).
Choose \(\varepsilon=K^{-2}\) in
\eqref{analytic:configuration-raw}. The result is the required
integrable version:
\begin{align}
 2\sum_{i<j}\log|t_i-t_j|-K\sum_iV(t_i)+\sum_i\sqrt{t_i}
 \leq{}&(\lambda M_0-I(\rho))K^2\notag\\
       &+(120+\sqrt2)h+2h\log K.
 \label{analytic:configuration}
\end{align}

\subsection{Restoring all scaling factors}

We now justify the reduction to \(V\), including an explicit error.
The power-series identity
\[
 \sum_{\ell\geq1}\ell^4q^\ell
       =\frac{q(1+11q+11q^2+q^3)}{(1-q)^5}
 \qquad(0<q<1)
\]
follows by applying \(q\,d/dq\) four times to \((1-q)^{-1}\).
Since \(1-e^{-u}\geq u/(1+u)\), substitution \(u=2\pi y\) in the
definition of \(w\) gives
\[
 w(y)\leq\frac1\pi(1+2\pi y)^5e^{-2\pi y}
       \leq8192(1+y)^5e^{-2\pi y}.
\]
For the last inequality use \(1+2\pi y\leq2\pi(1+y)\) and \(\pi<4\).

For \(1\leq m\leq K\), the right-endpoint Riemann sum satisfies
\begin{equation}\label{analytic:riemann}
 0\leq \sum_{j=1}^m\log(t+(j/K)^2)
               -K\int_0^{m/K}\log(t+u^2)\,du
       \leq2\log K+2,\qquad t\geq0.
\end{equation}
The lower bound is monotonicity in \(u\).
For \(t>0\), differentiating the difference gives the error of a
right-endpoint sum for the decreasing function \(1/(t+u^2)\);
that derivative is nonpositive. The largest difference is therefore
its limit at \(t=0\), namely
\[
 2\log(m!)-2m\log m+2m\leq2\log m+2\leq2\log K+2.
\]
Here \(\log(m!)\leq m\log m-m+\log m+1\), obtained by comparison
with \(\int_1^m\log u\,du\).
The continuous limit at zero follows either by integrating
\(\log(t+u^2)\) explicitly or by its integrable logarithmic limit.

Apply \eqref{analytic:riemann} to the numerator and denominator
products. Because the denominator logarithm is subtracted, only the
numerator's upper error is needed:
\[
 A_K(K^2t)\leq K^{12N-2K+12}e^{12}
  \exp\left(K\left[6\int_0^\alpha\log(t+u^2)\,du
                         -\int_0^1\log(t+u^2)\,du\right]\right).
\]
Under \(y_i=K\sqrt{t_i}\), each differential contributes
\(K/(2\sqrt{t_i})\), each weight contributes at most
\(8192K^5(1+\sqrt{t_i})^5e^{-2\pi K\sqrt{t_i}}\), and the
Vandermonde contributes \(K^{2h(h-1)}\).
Multiplying these factors in \eqref{analytic:andreief} gives
\begin{align}
 \Delta_K(\zeta(5))\leq{}&
 \frac{K^{2h(h+6N-K)+16h}(4096e^{12})^h}{h!}
 \int_{(0,\infty)^h}\prod_{i<j}(t_i-t_j)^2
 \prod_i e^{-KV(t_i)}t_i^{-1/2}(1+\sqrt{t_i})^5\,dt_i.
 \label{analytic:scaled-integral}
\end{align}
This keeps every power of \(K\) visible; the derivation is
\cite[(6.11)--(6.13)]{fauzan}.

By \eqref{analytic:configuration}, the remaining multiple integral is
at most
\[
 \exp\bigl((\lambda M_0-I(\rho))K^2
                      +(120+\sqrt2)h+2h\log K\bigr)\,652^h,
\]
because
\[
 \int_0^\infty t^{-1/2}(1+\sqrt t)^5e^{-\sqrt t}\,dt
 =2\int_0^\infty(1+x)^5e^{-x}\,dx
 =2\sum_{j=0}^5\binom5j j!=652.
\]
Dropping the favourable term \(-\log h!\), and using
\(\log4096<8.4,\ \log652<7,\ \sqrt2<1.5\), yields
\begin{equation}\label{analytic:delta-bound}
 \log\Delta_K(\zeta(5))
 \leq2h(h+6N-K)\log K+(\lambda M_0-I(\rho))K^2
                         +18h\log K+160h.
\end{equation}

\subsection{The factorial normalization and the final constant}

The normalization is
\[
 S_K=\frac{(K!)^{2h}4^{h-1}}
              {(N!)^{12h}\prod_{i=1}^{h-1}((2i)!)^2}.
\]
We give enough detail to check the cancellation of the large terms.
Integral comparisons imply, for every integer \(m\geq1\),
\[
 m\log m-m\leq\log(m!)
       \leq m\log m-m+\log m+1.
\]
Also
\begin{equation}\label{analytic:factorial-sum}
 \sum_{i=1}^{h-1}\log((2i)!)
       \geq h^2\log(2h)-\frac32h^2-2h\log(2h).
\end{equation}
For a direct proof put \(f(x)=2x\log(2x)-2x\), with \(f(0)=0\).
The factorial lower bound gives \(\log((2i)!)\geq f(i)\).
The function \(f\) increases for \(x\geq1/2\), so
\(\int_{i-1}^{i}f(x)\,dx\leq f(i)\) for \(i\geq2\).
For \(i=1\), direct integration gives the same inequality because
\(\log2\geq1/2\). Hence
\[
 \sum_{i=1}^{h-1}f(i)
 \geq\int_0^h f(x)\,dx-f(h)
 =h^2\log(2h)-\frac32h^2-2h\log(2h)+2h,
\]
which proves \eqref{analytic:factorial-sum}.

Substitute the three bounds into \(\log S_K\) and use \(N=\alpha K\),
\(h=\lambda K\). With
\begin{equation}\label{analytic:cstar}
 C_*=-2\lambda+12\alpha\lambda(1-\log\alpha)
                    +3\lambda^2-2\lambda^2\log(2\lambda),
\end{equation}
we obtain
\begin{align}
 \log S_K\leq{}&
 (2\lambda-12\alpha\lambda-2\lambda^2)K^2\log K+C_*K^2
                       +6h\log K+7h. \label{analytic:s-bound}
\end{align}
For the final \(7h\), the remaining constant is at most
\([2+4\log(2\lambda)+\log4]h\); use
\(\log(2\lambda)\leq2\lambda-1=0.85\) and \(\log4<1.4\).
This slightly relaxed linear error avoids any need for sharp Stirling
constants. The leading terms and normalization calculation are those of
\cite[(6.3), (6.15)]{fauzan}, also verified in
\cite[\texttt{Apery/RealBound.lean}, \texttt{log\_S\_le}]{firsching}.

The two coefficients of \(K^2\log K\) cancel exactly:
\[
 2\lambda(\lambda+6\alpha-1)
       +(2\lambda-12\alpha\lambda-2\lambda^2)=0.
\]
Thus \eqref{analytic:delta-bound} and \eqref{analytic:s-bound} imply
\begin{equation}\label{analytic:combined}
 \log F_K(\zeta(5))
 \leq(\lambda M_0-I(\rho)+C_*)K^2+24h\log K+167h.
\end{equation}
There is now just one numerical comparison left. Using
\eqref{analytic:energy}, \eqref{analytic:cstar}, and the logarithm
enclosures \eqref{analytic:log-series}, the same exact verifier gives
\[
 \begin{split}
 -2.126593445147050404
       &< I(\rho)<-2.126593445147050403,\\
  2.653035990340488661
       &< C_*<2.653035990340488662.
 \end{split}
\]
Using \(\lambda M_0=-49173/8000\), a direct interval computation of
the combination gives
\begin{equation}\label{analytic:constant}
 \lambda M_0-I(\rho)+C_*
 <-1.366995564512460935
 <-\frac{2733991}{2000000}=-1.3669955.
\end{equation}
These are outward rational bounds, not rounded guesses.
Fauzan gives this value and margin in Appendix A.4 \cite{fauzan}.
The independent interval check accompanies this draft.
The corresponding formal proof is the theorem
\texttt{energy\_const} in Firsching's \texttt{EnergyConst.lean} \cite{firsching}.
Finally \(h\leq K\), \(K\geq40\), and \(\log K\geq0\) turn
\eqref{analytic:combined} into \eqref{analytic:goal}.
The real estimate is therefore established with the same normalization
that appears in the arithmetic argument.
